\chapter{Patterns: data flow}
\label{ch:dataflow}

\begin{aside}
One-way data flow: state down, events up. The opposite leads
to tangle.
\end{aside}

\section{The pattern}

Parents pass data to children as props. Children emit
events that parents handle.

\begin{lstlisting}
auto editable = [&] (std::string val, auto on_change) {
    return text_input(val, on_change);
};

auto parent = [&] {
    auto v = signal("");
    return editable(ctx.get(v), [&](auto s) {
        ctx.set(v, s);
    });
};
\end{lstlisting}

\section{Two-way binding}

Some frameworks have two-way binding (the child can mutate
the parent's state directly). \maya{} doesn't. Pass a
mutator callback if needed:

\begin{lstlisting}
auto field = bound_input(signal_ref);
\end{lstlisting}

The helper bundles read and write into one parameter.

\section{Lifting state}

When two children need to share state, lift to their parent.
The shared signal lives in the parent; both children read
and write via parent-provided callbacks.

\section{Unidirectional reasoning}

The mental model: data flows down, events bubble up. State
lives at the highest level needed. Each component is a
function: same props in, same render out.

\section{Recap}

\begin{recap}
\begin{itemize}[leftmargin=*]
  \item Data down, events up.
  \item No two-way binding; explicit callbacks.
  \item Lift shared state to common parent.
\end{itemize}
\end{recap}

\section*{Workshop}

\begin{enumerate}[leftmargin=*]
  \item Identify any reverse flow in your app.
  \item Refactor to one-way.
\end{enumerate}
